\section{Bố cục luận văn}
Ngoài phần Tóm tắt, Danh mục tài liệu tham khảo, nội dung chính của luận văn được tổ chức thành 5 chương như sau:
\begin{itemize}
    \item Chương 1: Giới thiệu tổng quan đề tài. Chương này trình bày lý do chọn đề tài, tính cấp thiết của đề tài trong việc nghiên cứu các mô hình học sâu giải quyết bài toán tái tạo mô hình 3D chính xác cho vấn đề quản lý hạ tầng viễn thông. Đồng thời, chương xác định rõ mục tiêu nghiên cứu, đối tượng, phạm vi nghiên cứu, phương pháp nghiên cứu và những đóng góp nghiên cứu khoa học chính của luận văn. Đồng thời trình bày bố cục tổng thể của các chương tiếp theo.

    \item Chương 2: Cơ sở lý thuyết. Chương này cung cấp nền tảng kiến thức lý thuyết cần thiết cho luận văn. Nội dung bao gồm giới thiệu về bài toán tái tạo đám mây điểm 3D và chi tiết, cách tiếp cận ước tính camera pose của thuật toán SfM truyền thống và các thuật toán trích xuất đặc trưng hiện đại. Bên cạnh đó, chương này cũng giới thiệu các phương pháp tái tạo đám mây điểm 3D chi tiết, dày đặc phổ biến: MVS và NeRF.

    \item Chương 3: Phương pháp đề xuất. Chương này tập trung trình bày chi tiết kỹ thuật đề xuất Hybrid-SfM nhằm nâng cao độ chính xác khi tính toán camera pose cho dữ liệu ảnh cột ăng-ten, trạm viễn thông. 
    
    \item Chương 4: Thực nghiệm và thảo luận. Chương này trình bày quá trình thiết lập và tiến hành các thực nghiệm nhằm kiểm chứng và đánh giá hiệu quả của kỹ thuật đề xuất. Nội dung bao gồm mô tả chi tiết về các bộ dữ liệu được sử dụng, các độ đo đánh giá, môi trường cài đặt và các kịch bản thực nghiệm. Các kết quả định lượng và định tính thu được sẽ được trình bày, so sánh giữa các phương pháp với nhau. Cuối chương là phần thảo luận, phân tích sâu về kết quả, chỉ ra ưu nhược điểm của phương pháp đề xuất khi áp dụng vào bài toán thực tế là tái tạo mô hình đám mây điểm 3D chất lượng cao cho cột ăng-ten của trạm viễn thông.
    
    \item Chương 5: Kết luận. Phần cuối cùng tóm tắt lại những kết quả chính và đóng góp cốt lõi của luận văn. Đồng thời, chương cũng chỉ ra những hạn chế còn tồn tại của nghiên cứu và đề xuất các hướng phát triển, cải tiến tiềm năng trong tương lai.
\end{itemize}

